%------------------------------------------------------------------------------%
\documentclass[a4paper,12pt,fleqn]{article}

\usepackage{gaal}

\ShowAnswers

%------------------------------------------------------------------------------%
\begin{document}

\makeHeader{GAAL}{Prova 1 -- S}{2026-2}

% 01
%------------------------------------------------------------------------------%
\exercise{25\%}
Dados
\[
A=
\begin{bmatrix}
  1 &  2 & -1\\
  0 &  1 &  3
\end{bmatrix}
\qquad
B=
\begin{bmatrix}
  2 & -1 &  0\\
 -1 &  2 &  1
\end{bmatrix}
\qquad
C=
\begin{bmatrix}
  1 &  0 &  2\\
  2 & -1 &  1
\end{bmatrix}
\]
calcule
\[
(A+B)C^t+C(B^t+A^t)
\]

\begin{answer}
Utilizamos a propriedade da transposta de uma soma
\[
B^t+A^t=(A+B)^t
\]
Assim
\begin{align*}
&(A+B)C^t+C(B^t+A^t)
\\
&=(A+B)C^t+C(A+B)^t
\end{align*}
Primeiramente
\begin{align*}
A+B
&=
\begin{bmatrix}
  1 &  2 & -1\\
  0 &  1 &  3
\end{bmatrix}
+
\begin{bmatrix}
  2 & -1 &  0\\
 -1 &  2 &  1
\end{bmatrix}
\\[2mm]
&=
\begin{bmatrix}
  3 & 1 & -1\\
 -1 & 3 &  4
\end{bmatrix}
\end{align*}
Calculamos o primeiro produto
\begin{align*}
(A+B)C^t
&=
\begin{bmatrix}
  3 & 1 & -1\\
 -1 & 3 &  4
\end{bmatrix}
\begin{bmatrix}
  1 & 2\\
  0 & -1\\
  2 & 1
\end{bmatrix}
\\[2mm]
&=
\begin{bmatrix}
  3-2 & 6-1-1\\
 -1+8 & -2-3+4
\end{bmatrix}
\\[2mm]
&=
\begin{bmatrix}
  1 & 4\\
  7 & -1
\end{bmatrix}
\end{align*}
Pela propriedade da transposta de um produto
\[
\left[(A+B)C^t\right]^t=C(A+B)^t
\]
Portanto
\[
C(B^t+A^t)=
\begin{bmatrix}
  1 & 7\\
  4 & -1
\end{bmatrix}
\]
Somando as matrizes
\begin{align*}
&(A+B)C^t+C(B^t+A^t)
\\[2mm]
&=
\begin{bmatrix}
  1 & 4\\
  7 & -1
\end{bmatrix}
+
\begin{bmatrix}
  1 & 7\\
  4 & -1
\end{bmatrix}
\\[2mm]
&=
\begin{bmatrix}
  2 & 11\\
  11 & -2
\end{bmatrix}
\end{align*}
Portanto
\[
\boxed{
\begin{bmatrix}
  2 & 11\\
  11 & -2
\end{bmatrix}
}
\]
Observe que o resultado é uma matriz simétrica.
\end{answer}

% 02
%------------------------------------------------------------------------------%
\exercise{25\%}
Resolva o sistema linear por Gauss--Jordan
\[
\systeme[sort={x,y,z,w}]{
  2x+y-z=9,
  x+y+z+4w=4,
  -x+2y+z-w=0
}
\]

\begin{answer}
A matriz aumentada do sistema é
\[
\begin{amatrix}{4}
  2 & 1 & -1 &  0 & 9\\
  1 & 1 &  1 &  4 & 4\\
 -1 & 2 &  1 & -1 & 0
\end{amatrix}
\]
Trocamos a primeira linha pela segunda
\[
L_1\leftrightarrow L_2
\]
Obtemos
\[
\begin{amatrix}{4}
  1 & 1 &  1 &  4 & 4\\
  2 & 1 & -1 &  0 & 9\\
 -1 & 2 &  1 & -1 & 0
\end{amatrix}
\]
Utilizamos a primeira linha para eliminar os elementos abaixo do primeiro pivô
\[
L_2\leftarrow L_2-2L_1
\]
\[
L_3\leftarrow L_3+L_1
\]
Assim
\[
\begin{amatrix}{4}
  1 & 1  &  1 &  4 & 4\\
  0 & -1 & -3 & -8 & 1\\
  0 & 3  &  2 &  3 & 4
\end{amatrix}
\]
Multiplicamos a segunda linha por $-1$
\[
L_2\leftarrow -L_2
\]
Logo
\[
\begin{amatrix}{4}
  1 & 1 & 1 & 4 & 4\\
  0 & 1 & 3 & 8 & -1\\
  0 & 3 & 2 & 3 & 4
\end{amatrix}
\]
Eliminamos o elemento abaixo do segundo pivô
\[
L_3\leftarrow L_3-3L_2
\]
Obtemos
\[
\begin{amatrix}{4}
  1 & 1 &  1 &  4 &  4\\
  0 & 1 &  3 &  8 & -1\\
  0 & 0 & -7 & -21 & 7
\end{amatrix}
\]
Dividimos a terceira linha por $-7$
\[
L_3\leftarrow -\frac{1}{7}L_3
\]
Assim
\[
\begin{amatrix}{4}
  1 & 1 & 1 & 4 &  4\\
  0 & 1 & 3 & 8 & -1\\
  0 & 0 & 1 & 3 & -1
\end{amatrix}
\]
Eliminamos os elementos acima do terceiro pivô
\[
L_2\leftarrow L_2-3L_3
\]
\[
L_1\leftarrow L_1-L_3
\]
Resulta
\[
\begin{amatrix}{4}
  1 & 1 & 0 &  1 & 5\\
  0 & 1 & 0 & -1 & 2\\
  0 & 0 & 1 &  3 & -1
\end{amatrix}
\]
Finalmente eliminamos o elemento acima do segundo pivô
\[
L_1\leftarrow L_1-L_2
\]
Obtemos
\[
\begin{amatrix}{4}
  1 & 0 & 0 &  2 &  3\\
  0 & 1 & 0 & -1 &  2\\
  0 & 0 & 1 &  3 & -1
\end{amatrix}
\]
A matriz escalonada reduzida representa o sistema
\[
\systeme[sort={x,y,z,w}]{
  x+2w=3,
  y-w=2,
  z+3w=-1
}
\]
A variável $w$ é livre
\[
w=t
\qquad t\in\R
\]
Então
\begin{align*}
x&=3-2t\\
y&=2+t\\
z&=-1-3t
\end{align*}
Portanto, a solução geral é
\[
\begin{pmatrix}
x\\y\\z\\w
\end{pmatrix}
=
\begin{pmatrix}
3-2t\\
2+t\\
-1-3t\\
t
\end{pmatrix}
\]
ou
\[
\boxed{
\begin{pmatrix}
x\\y\\z\\w
\end{pmatrix}
=
\begin{pmatrix}
3\\2\\-1\\0
\end{pmatrix}
+
t
\begin{pmatrix}
-2\\1\\-3\\1
\end{pmatrix}
\qquad t\in\R
}
\]
\end{answer}

% 03
%------------------------------------------------------------------------------%
\exercise{25\%}
Utilize o método de Gauss--Jordan para inverter a matriz
\[
A=
\begin{bmatrix}
  1 & 2 & 1\\
  2 & 5 & 3\\
  1 & 1 & 2
\end{bmatrix}
\]

\begin{answer}
Construímos a matriz aumentada
\[
\begin{amatrix}[3]{3}
  1 & 2 & 1 & 1 & 0 & 0\\
  2 & 5 & 3 & 0 & 1 & 0\\
  1 & 1 & 2 & 0 & 0 & 1
\end{amatrix}
\]
Eliminamos os elementos abaixo do primeiro pivô
\[
L_2\leftarrow L_2-2L_1
\]
\[
L_3\leftarrow L_3-L_1
\]
Obtemos
\[
\begin{amatrix}[3]{3}
  1 &  2 &  1 &  1 &  0 &  0\\
  0 &  1 &  1 & -2 &  1 &  0\\
  0 & -1 &  1 & -1 &  0 &  1
\end{amatrix}
\]
Eliminamos o elemento abaixo do segundo pivô
\[
L_3\leftarrow L_3+L_2
\]
Assim
\[
\begin{amatrix}[3]{3}
  1 & 2 & 1 &  1 & 0 & 0\\
  0 & 1 & 1 & -2 & 1 & 0\\
  0 & 0 & 2 & -3 & 1 & 1
\end{amatrix}
\]
Dividimos a terceira linha por $2$
\[
L_3\leftarrow\frac{1}{2}L_3
\]
Logo
\[
\begin{amatrix}[3]{3}
  1 & 2 & 1 &  1 & 0 & 0\\
  0 & 1 & 1 & -2 & 1 & 0\\
  0 & 0 & 1 & -\frac{3}{2} & \frac{1}{2} & \frac{1}{2}
\end{amatrix}
\]
Eliminamos os elementos acima do terceiro pivô
\[
L_2\leftarrow L_2-L_3
\]
\[
L_1\leftarrow L_1-L_3
\]
Obtemos
\[
\begin{amatrix}[3]{3}
  1 & 2 & 0 & \frac{5}{2} & -\frac{1}{2} & -\frac{1}{2}\\
  0 & 1 & 0 & -\frac{1}{2} & \frac{1}{2} & -\frac{1}{2}\\
  0 & 0 & 1 & -\frac{3}{2} & \frac{1}{2} & \frac{1}{2}
\end{amatrix}
\]
Finalmente eliminamos o elemento acima do segundo pivô
\[
L_1\leftarrow L_1-2L_2
\]
Resulta
\[
\begin{amatrix}[3]{3}
  1 & 0 & 0 & \frac{7}{2} & -\frac{3}{2} & \frac{1}{2}\\
  0 & 1 & 0 & -\frac{1}{2} & \frac{1}{2} & -\frac{1}{2}\\
  0 & 0 & 1 & -\frac{3}{2} & \frac{1}{2} & \frac{1}{2}
\end{amatrix}
\]
Portanto
\[
\boxed{
A^{-1}=
\begin{bmatrix}
 \frac{7}{2} & -\frac{3}{2} &  \frac{1}{2}\\
-\frac{1}{2} &  \frac{1}{2} & -\frac{1}{2}\\
-\frac{3}{2} &  \frac{1}{2} &  \frac{1}{2}
\end{bmatrix}
}
\]
\end{answer}

% 04
%------------------------------------------------------------------------------%
\exercise{25\%}
Utilize a inversa para resolver o sistema linear
\[
\systeme{
  x+2y+z=7,
  2x+5y+3z=13,
  x+y+2z=8
}
\]
\hint{Utilize a matriz inversa calculada no exercício anterior}

\begin{answer}
A matriz do sistema é
\[
A=
\begin{bmatrix}
  1 & 2 & 1\\
  2 & 5 & 3\\
  1 & 1 & 2
\end{bmatrix}
\]
e o vetor dos termos independentes é
\[
\mathbf{b}=
\begin{bmatrix}
7\\13\\8
\end{bmatrix}
\]
Do exercício anterior temos
\[
A^{-1}=
\begin{bmatrix}
 \frac{7}{2} & -\frac{3}{2} &  \frac{1}{2}\\
-\frac{1}{2} &  \frac{1}{2} & -\frac{1}{2}\\
-\frac{3}{2} &  \frac{1}{2} &  \frac{1}{2}
\end{bmatrix}
\]
A solução do sistema é
\begin{align*}
\mathbf{x}
&=A^{-1}\mathbf{b}
\\[2mm]
&=
\begin{bmatrix}
 \frac{7}{2} & -\frac{3}{2} &  \frac{1}{2}\\
-\frac{1}{2} &  \frac{1}{2} & -\frac{1}{2}\\
-\frac{3}{2} &  \frac{1}{2} &  \frac{1}{2}
\end{bmatrix}
\begin{bmatrix}
7\\13\\8
\end{bmatrix}
\\[2mm]
&=
\begin{bmatrix}
\frac{49}{2}-\frac{39}{2}+\frac{8}{2}\\
-\frac{7}{2}+\frac{13}{2}-\frac{8}{2}\\
-\frac{21}{2}+\frac{13}{2}+\frac{8}{2}
\end{bmatrix}
\\[2mm]
&=
\begin{bmatrix}
\frac{18}{2}\\
-\frac{2}{2}\\
0
\end{bmatrix}
\\[2mm]
&=
\begin{bmatrix}
9\\-1\\0
\end{bmatrix}
\end{align*}
Portanto
\[
\boxed{
(x,y,z)=(9,-1,0)
}
\]
\end{answer}

%------------------------------------------------------------------------------%
\end{document}
%------------------------------------------------------------------------------%
